% 本文件由 LaTeX 排版 Agent 根据有效 translation.json 自动生成。
% author=Ignatius of Antioch; work_id=0123; title=依纳爵殉道记

\chapter{依纳爵殉道记}

% structure: p0001 source=h1 target=omitted reason=page-title-already-rendered-by-parent

% structure: p0002 source=h2 target=section reason=relative-heading-rank; base=h2

\section{第一章 依纳爵渴望殉道}

不久以前，图拉真继承了罗马帝国的帝位，宗徒若望的门徒依纳爵，一个在各方面具有宗徒品格的人，以极大的谨慎治理安提约基雅教会；他在多米提安治下多次迫害的暴风雨中艰难逃脱，如同一位优秀的舵手，藉着祈祷和斋戒的船舵，藉着他教导的热忱，以及他不断的属灵劳作，抵挡了扑面而来的洪流，他惟恐失去那些缺乏勇气或因其纯朴而易于受苦的人。因此，当迫害暂时止息，教会处于平安时，他感到喜乐，但为自己感到忧伤，因为他尚未达到对基督真正的爱，也未达到门徒的完美等级。因为他内心反思，藉殉道所作的宣认，会使他与主有更亲密的关系。因此，他继续与教会共处数年，如同神圣的明灯，藉着他对圣经的阐释光照每个人的心智，终于达到了他的愿望。\par

% structure: p0004 source=h2 target=section reason=relative-heading-rank; base=h2

\section{第二章 依纳爵被图拉真定罪}

因为图拉真在他统治的第九年，在对斯基泰人、达契亚人以及许多其他民族取得胜利后，心高气傲，认为基督信徒的宗教团体是他完全征服万物的障碍，于是威胁他们，若不敬拜鬼神如同其他民族一样，就加以迫害，从而迫使所有过虔诚生活的人要么向偶像献祭，要么死亡。因此，基督的卓越战士依纳爵，为安提约基雅教会担忧，按照他自己的愿望，被带到当时驻在安提约基雅、但正急于出发征讨亚美尼亚和帕提亚人的图拉真面前。当他被带到皇帝图拉真面前时，皇帝对他说：\Quote{你是什么人，邪恶的坏蛋，竟敢违背我们的命令，还说服别人也这样做，使他们悲惨地灭亡？}依纳爵回答说：\Quote{没有人应该称特奥福鲁斯为邪恶的，因为一切邪灵都已从天主的仆人身上离去。但如果因为我是这些邪灵的敌人，你在这一点上称我为邪恶，我完全同意你的说法；因为我把天上的君王基督怀在心中，我摧毁了这些邪灵的一切诡计。}图拉真回答说：\Quote{特奥福鲁斯是谁？}依纳爵说：\Quote{那把基督怀在胸中的人。}图拉真说：\Quote{难道\Emph{我们}不是把诸神放在我们心中吗？我们与仇敌作战时得到他们的帮助。}依纳爵回答说：\Quote{你错了，你把民族的鬼神称作神。因为只有一位天主，他创造了天、地、海和其中的万物；还有一位耶稣基督，天主的独生子，我愿得享他的国。}图拉真说：\Quote{你是指那个在本丢·比拉多治下被钉死的人吗？}依纳爵回答说：\Quote{我是指那钉死我的罪和罪的创始者，并把魔鬼的一切欺骗和恶意都踏在那些心中怀着祂之人脚下的人。}图拉真说：\Quote{难道你里面怀着被钉死的那一位吗？}依纳爵说：\Quote{确实如此；因为经上记着：\InnerQuote{我要在他们内居住，并在他们中行走。}}\BibleRef{格后 6:16}于是图拉真宣判如下：\Quote{我们命令，那声称自己里边怀着被钉死之人的依纳爵，要被士兵捆绑，带到大城罗马，在那里被野兽吞噬，以飨民众。}圣殉道者听到这判决，喜乐地喊道：\Quote{主，我感谢你，因为你屈尊以对你的完美之爱来荣耀我，并使我像你的宗徒保禄一样被铁链捆绑。}说完这话，他欢喜地拥抱锁链；他首先为教会祈祷，流着泪将教会托付给主，然后被士兵残暴地拖走，如同一头卓越的公羊，一群美善羊群的领袖，被带往罗马，在那里成为嗜血野兽的食物。\par

% structure: p0006 source=h2 target=section reason=relative-heading-rank; base=h2

\section{第三章 依纳爵航行至斯米纳}

因此，带着极大的热忱和喜乐，出于对受苦的渴望，他从安提约基雅下到塞琉西亚，在那里乘船出发。经历许多苦难后，他到达斯米纳，带着极大的喜乐上岸，急忙去见圣波利卡普，他昔日的同门，当时是斯米纳的主教。因为他们两人在古时都曾是圣宗徒若望的门徒。他们相见后，依纳爵分给了他一些属神恩赐，并以自己的锁链为荣，请求他与他一同劳苦，以达成他的愿望；他热切地恳求整个教会（因为亚洲的城邑和教会都通过他们的主教、长老和执事欢迎了这位圣人，众人都急忙来见他，希望能从他那里得到某种属灵恩赐），尤其是圣波利卡普，使他能藉着野兽，很快离开这个世界，在基督面前显现。\par

% structure: p0008 source=h2 target=section reason=relative-heading-rank; base=h2

\section{第四章 依纳爵致信各教会}

他说这些话，作这见证，将他对基督的爱延伸到极致，如同一个将要通过美好的宣认以及那些为他的即将来临的战斗而联合祈祷之人的热忱而获得天堂的人；并回报那些通过他们的领袖来迎接他的教会，写感谢信给他们，信函滴下属神的恩宠，连同祈祷和劝勉。因此，看到所有人都对他如此友善，惟恐弟兄们的爱心会妨碍他对主的热情，正当受苦殉道的敞开之门摆在他面前时，他写信给罗马教会，即下面所附的那封信。\par

% structure: p0010 source=h2 target=section reason=relative-heading-rank; base=h2

\section{第五章 依纳爵被带到罗马}

因此，藉着这封信，他如愿安顿了罗马弟兄中那些不愿他殉道的人；他从斯米纳起航（因为士兵催促背负基督者\OriginalTerm{背负基督者}{Christophorus}赶紧前往大城罗马的公开竞技，好把他丢给野兽，在罗马民众眼前，使他获得他所追求的冠冕），随后他在特洛阿登陆。然后从那里前往纳玻里，徒步经过斐理伯，穿过马其顿，直到厄毗鲁斯靠近厄毗达姆努斯的地方；在一个海港找到船后，他航过亚得里亚海，由此进入第勒尼安海，经过各个岛屿和城市，直到普特奥利映入眼帘时，他急切想在那里上岸，渴望追随宗徒保禄的足迹。\BibleRef{宗 28:13}但暴风骤起，不容他这样做，船被迅速吹向前方；他只能对那地方弟兄们的爱德表示欣喜，便驶过了。因此，顺风持续，我们被迫在一天一夜内匆匆前行，为这位义人即将离开我们而哀伤。但对他而言，这正合他意，因为他急于尽快离开此世，好能到达他所爱的主那里。随后我们驶入罗马港口，不圣洁的竞技正要结束时，士兵们开始对我们的迟缓感到恼怒，而这位主教却欣喜地顺从了他们的催促。\par

% structure: p0012 source=h2 target=section reason=relative-heading-rank; base=h2

\section{第六章 依纳爵在罗马被野兽吞噬}

于是他们从名为波图斯的地方出发；有关圣殉道者的一切传闻已经传开，我们遇到充满恐惧与喜乐的弟兄们：他们因被堪当迎接背负天主者\OriginalTerm{背负天主者}{Theophorus}而喜乐，却又因如此杰出的人被引向死亡而恐惧。他吩咐一些因热心激昂而说要安抚民众、使他们不要求消灭这位义人的人保持沉默。他藉着圣神立即察觉此事，便问候了众人，恳求他们对他显示出真正的爱，并在这一点上比在他的书信中更详细地讲论，说服他们不要嫉妒他奔向主。然后，在全体弟兄跪在他身旁时，他为教会恳求天主圣子，使迫害得以停止，弟兄间的彼此相爱得以延续。之后，他便被急忙领进竞技场。随即，按照凯撒早些时候发出的命令，他被抛了进去，此时公开竞技正要结束（因为那天在他们看来是隆重之日，即罗马语中称为\Quote{第十三日}的那天，民众惯常聚集得比平时更多）。他被丢在靠近神庙的野兽中间，这样，圣殉道者依纳爵的愿望便藉着它们而实现，正如经上所载：\Quote{义人的渴望是蒙天主悦纳的}\BibleRef{箴 10:24}，以致他不因遗骸的收集而给任何弟兄带来麻烦，正如他在书信中预先表达过希望这样结局的愿望。因为只有他圣骸的较坚硬部分被留下，这些部分被送到安提约基雅，用细麻布包裹，作为无价之宝留给圣教会，这恩宠出自殉道者。\par

% structure: p0014 source=h2 target=section reason=relative-heading-rank; base=h2

\section{第七章 依纳爵死后显现异象}

这些事发生在罗马历正月前的第十三日，即十二月二十日，当时苏拉和塞内齐奥第二次担任罗马执政官。我们亲眼目睹了这些事，整夜在屋内流泪，屈膝恳切祈祷，求主赐给我们软弱之人对所发生之事的充分确信。我们小睡片刻时，有些人看见真福依纳爵突然站在我们身旁，拥抱我们；另一些人看见他再次为我们祈祷；还有人看见他汗流浃背，仿佛刚从大劳苦中出来，站在主身旁。我们怀着极大的喜乐见证这些事，比较各自的异象后，我们颂赞万善之源的天主，并表达对圣殉道者福分的感受。现在我们将日子和时间告知你们，好使你们按他殉道的时间聚集，与基督的斗士和高贵殉道者共融，他践踏了魔鬼，并完成了因爱基督而渴望的赛程，在我们的主基督耶稣内；愿光荣和权能藉着祂并同着祂，归于父及圣神，直到永远！阿们。\par

\SourceRecord{Translated by Alexander Roberts and James Donaldson.；Ante-Nicene Fathers；Vol. 1.；Edited by Alexander Roberts, James Donaldson, and A. Cleveland Coxe.；Buffalo, NY: Christian Literature Publishing Co.,；1885.；<http://www.newadvent.org/fathers/0123.htm>.}
